\documentclass[11pt,letterpaper]{article}
\usepackage{../../Common/course_notes}
\usepackage{../../Common/course_macros}

\begin{document}
\setlength{\abovedisplayskip}{6pt}
\setlength{\belowdisplayskip}{6pt}
\setlength{\textfloatsep}{10pt}
\setlength{\intextsep}{8pt}
\color{blackmain}

\lectureheader{MÓDULO 00}{SEMANA 04}{Matrices Jacobianas, Coordenadas Curvilíneas y el Tensor Métrico}{Bases locales adaptadas, factores de escala de Lamé y elementos de volumen}

% =====================================================================
\manualsection[1. De Coordenadas Cartesianas a Curvilíneas]{1. La Necesidad de Nuevas Coordenadas: Simetría Física frente a Rigidez Cartesiana}

En la física y la geometría aplicada rara vez nos enfrentamos a problemas donde la simetría sea estrictamente cúbica o plana. Cuando estudiamos el campo gravitatorio de un planeta, la difusión térmica en un cilindro o la rotación de un vórtice, forzar coordenadas cartesianas $(x, y, z)$ complica innecesariamente las ecuaciones.

Sin embargo, al abandonar el sistema cartesiano surge un desafío matemático inmediato: las líneas coordenadas ya no son rectas paralelas y los vectores unitarios cambian de dirección de un punto a otro del espacio. ¿Cómo se transforman los vectores, los desplazamientos infinitesimales y las distancias en un sistema de coordenadas no cartesiano?

% =====================================================================
\manualsection[2. La Matriz Jacobiana y la Base Local]{2. La Matriz Jacobiana y la Construcción de la Base Local}

Consideremos una transformación suave de coordenadas en $\R^n$ entre variables cartesianas $x^i$ ($i = 1, \dots, n$) y coordenadas curvilíneas generales $u^a$ ($a = 1, \dots, n$):
\[
x^i = x^i(u^1, u^2, \dots, u^n).
\]

\begin{intuitionbox}{El Desplazamiento Infinitesimal y la Matriz Jacobiana}
Por la regla de la cadena multivariable, la variación infinitesimal de las coordenadas cartesianas al variar las coordenadas curvilíneas está dada por:
\[
dx^i = \frac{\partial x^i}{\partial u^a}\,du^a \equiv J^i{}_a\,du^a,
\]
donde hemos adoptado la convención de suma de Einstein sobre el índice mudo $a$.
La matriz $J = (J^i{}_a)_{n\times n}$ se denomina \textbf{Matriz Jacobiana}:
\[
J = \begin{pmatrix} 
\frac{\partial x^1}{\partial u^1} & \cdots & \frac{\partial x^1}{\partial u^n} \\
\vdots & \ddots & \vdots \\
\frac{\partial x^n}{\partial u^1} & \cdots & \frac{\partial x^n}{\partial u^n}
\end{pmatrix}.
\]
\end{intuitionbox}

\begin{theorembox}{Construcción de la Base Local Tangente}
El vector de posición cartesiano es $\vec{r}(u^1, \dots, u^n) = x^i(u^a)\hat{e}_i$. Si incrementamos una sola coordenada $u^a$ manteniendo todas las demás fijas, la tasa de variación espacial genera el vector tangente a la $a$-ésima línea coordenada:
\[
\vec{\varepsilon}_a \equiv \frac{\partial\vec{r}}{\partial u^a} = \frac{\partial x^i}{\partial u^a}\hat{e}_i = J^i{}_a \hat{e}_i.
\]
El conjunto $\{\vec{\varepsilon}_1, \dots, \vec{\varepsilon}_n\}$ constituye la \textbf{base local (o coordenada)} en el punto $P$. A diferencia de la base cartesiana fija $\{\hat{e}_i\}$, los vectores $\vec{\varepsilon}_a$ dependen del punto donde nos encontremos.
\end{theorembox}

\begin{center}
\begin{tikzpicture}[scale=1.0]
    % Ejes cartesianos suaves
    \draw[-{Stealth}, color=graymuted] (-0.3,0) -- (4.2,0) node[right, font=\footnotesize] {$x$};
    \draw[-{Stealth}, color=graymuted] (0,-0.3) -- (0,3.6) node[above, font=\footnotesize] {$y$};

    \coordinate (O) at (0,0);
    \coordinate (P) at (2.4, 1.8);

    % Linea coordenada r (radial)
    \draw[color=orangeacc!70, domain=0.4:3.6, variable=\t, samples=20] plot ({\t}, {0.75*\t});
    \node[color=orangeacc, font=\scriptsize\bfseries] at (3.5, 2.8) {Línea $r$ ($\theta = \text{cte}$)};

    % Linea coordenada theta (arco)
    \draw[color=navyaccent!70, domain=20:65, variable=\a, samples=30] plot ({3.0*cos(\a)}, {3.0*sin(\a)});
    \node[color=navyaccent, font=\scriptsize\bfseries] at (1.1, 3.1) {Línea $\theta$ ($r = \text{cte}$)};

    % Vectores de la base local en P
    \draw[-{Stealth[length=2.8mm]}, line width=1.8pt, color=orangeacc] (P) -- ($(P) + (1.2, 0.9)$) node[above right=1pt, font=\small\bfseries] {$\vec{\varepsilon}_r = \frac{\partial\vec{r}}{\partial r}$};
    \draw[-{Stealth[length=2.8mm]}, line width=1.8pt, color=navyaccent] (P) -- ($(P) + (-0.9, 1.2)$) node[above left=1pt, font=\small\bfseries] {$\vec{\varepsilon}_\theta = \frac{\partial\vec{r}}{\partial \theta}$};

    \fill[blackmain] (P) circle (2.2pt) node[below right=2pt, font=\footnotesize\bfseries] {$P(r,\theta)$};
    \draw[-{Stealth}, color=graymuted, line width=1pt] (O) -- (P) node[pos=0.5, below right=1pt, font=\footnotesize] {$\vec{r}$};
\end{tikzpicture}
\end{center}

% =====================================================================
\manualsection[3. El Tensor Métrico]{3. El Tensor Métrico $g_{ab}$: Midiendo Distancias en Coordenadas Curvas}

En coordenadas cartesianas, el Teorema de Pitágoras diferencial establece que el cuadrado de la distancia infinitesimal entre dos puntos cercanos es:
\[
ds^2 = (dx^1)^2 + (dx^2)^2 + \dots + (dx^n)^2 = \delta_{ij} dx^i dx^j.
\]
Si sustituimos la transformación $dx^i = J^i{}_a du^a$ en esta expresión:
\[
ds^2 = \delta_{ij} (J^i{}_a du^a)(J^j{}_b du^b) = (\delta_{ij} J^i{}_a J^j{}_b) du^a du^b.
\]

\begin{theorembox}{Definición Formal: Tensor Métrico Riemanniano $g_{ab}$}
El objeto bilineal simétrico que gobierna el cálculo de distancias y productos internos en las coordenadas $u^a$ es el \textbf{Tensor Métrico}:
\[
\boxed{g_{ab} \equiv \vec{\varepsilon}_a \cdot \vec{\varepsilon}_b = \delta_{ij} J^i{}_a J^j{}_b = (J^T J)_{ab}}
\]
La distancia euclidiana al cuadrado (o elemento de línea) se expresa mediante la forma cuadrática fundamental:
\[
\boxed{ds^2 = g_{ab}\,du^a du^b}
\]
\end{theorembox}

\manualsubsection{Sistemas Coordenados Ortogonales}
Un sistema de coordenadas curvilíneas es \textbf{ortogonal} si sus vectores coordenados locales son mutuamente perpendiculares para todo punto del espacio:
\[
\vec{\varepsilon}_a \cdot \vec{\varepsilon}_b = 0 \quad \text{para todo } a \neq b.
\]
En tal caso, la matriz del tensor métrico es diagonal:
\[
(g_{ab}) = \begin{pmatrix} h_1^2 & 0 & \dots & 0 \\ 0 & h_2^2 & \dots & 0 \\ \vdots & \vdots & \ddots & \vdots \\ 0 & 0 & \dots & h_n^2 \end{pmatrix},
\]
donde $h_a \equiv \|\vec{\varepsilon}_a\| = \sqrt{g_{aa}}$ se denominan \textbf{factores de escala de Lamé}.

% =====================================================================
\manualsection[4. Coordenadas Clásicas: Polares, Cilíndricas y Esféricas]{4. Los Tres Ejemplos Arquetípicos de la Física}

\manualsubsection{1. Coordenadas Polares en el Plano $(r, \theta)$}
Las ecuaciones de transformación son $x = r\cos\theta$ e $y = r\sin\theta$.
La matriz Jacobiana es:
\[
J = \begin{pmatrix} \frac{\partial x}{\partial r} & \frac{\partial x}{\partial \theta} \\ \frac{\partial y}{\partial r} & \frac{\partial y}{\partial \theta} \end{pmatrix} = \begin{pmatrix} \cos\theta & -r\sin\theta \\ \sin\theta & r\cos\theta \end{pmatrix}.
\]
El tensor métrico se obtiene multiplicando $J^T J$:
\[
(g_{ab}) = J^T J = \begin{pmatrix} \cos\theta & \sin\theta \\ -r\sin\theta & r\cos\theta \end{pmatrix} \begin{pmatrix} \cos\theta & -r\sin\theta \\ \sin\theta & r\cos\theta \end{pmatrix} = \begin{pmatrix} 1 & 0 \\ 0 & r^2 \end{pmatrix}.
\]
Por lo tanto, el elemento de línea es:
\[
\boxed{ds^2 = dr^2 + r^2 d\theta^2}
\]
Los vectores tangentes tienen normas $\|\vec{\varepsilon}_r\| = 1$ y $\|\vec{\varepsilon}_\theta\| = r$. Los vectores unitarios normalizados son $\hat{e}_r = \vec{\varepsilon}_r$ y $\hat{e}_\theta = \frac{1}{r}\vec{\varepsilon}_\theta$.

\manualsubsection{2. Coordenadas Cilíndricas en $\R^3$ $(\rho, \phi, z)$}
Con $x = \rho\cos\phi$, $y = \rho\sin\phi$ y $z = z$:
\[
(g_{ab}) = \begin{pmatrix} 1 & 0 & 0 \\ 0 & \rho^2 & 0 \\ 0 & 0 & 1 \end{pmatrix} \quad \implies \quad \boxed{ds^2 = d\rho^2 + \rho^2 d\phi^2 + dz^2}
\]

\manualsubsection{3. Coordenadas Esféricas en $\R^3$ $(r, \theta, \phi)$}
Con $x = r\sin\theta\cos\phi$, $y = r\sin\theta\sin\phi$, $z = r\cos\theta$:
Los factores de escala son $h_r = 1$, $h_\theta = r$, $h_\phi = r\sin\theta$, lo que conduce a:
\[
(g_{ab}) = \begin{pmatrix} 1 & 0 & 0 \\ 0 & r^2 & 0 \\ 0 & 0 & r^2\sin^2\theta \end{pmatrix} \quad \implies \quad \boxed{ds^2 = dr^2 + r^2 d\theta^2 + r^2\sin^2\theta\,d\phi^2}
\]

% =====================================================================
\manualsection[5. Elementos de Volumen y el Determinante Jacobiano]{5. La Deformación del Volumen y el Determinante Jacobiano}

¿Cómo cambia un elemento de volumen cartesiano $dV = dx\,dy\,dz$ al integrarse en coordenadas curvilíneas?
El paralelepípedo infinitesimal generado por los vectores de la base local $\{\vec{\varepsilon}_1 du^1, \dots, \vec{\varepsilon}_n du^n\}$ tiene volumen:
\[
dV = |(\vec{\varepsilon}_1 \times \vec{\varepsilon}_2)\cdot\vec{\varepsilon}_3|\,du^1 du^2 du^3 = |\det(J)|\,du^1 du^2 du^3.
\]
Usando la relación métrica $g = J^T J$:
\[
\det(g) = \det(J^T J) = [\det(J)]^2 \quad \implies \quad |\det(J)| = \sqrt{\det(g)}.
\]
Por lo tanto:
\[
\boxed{dV = \sqrt{\det(g)}\,du^1 \dots du^n}
\]
Para coordenadas esféricas, $\det(g) = (1)(r^2)(r^2\sin^2\theta) = r^4\sin^2\theta$, de modo que $\sqrt{\det(g)} = r^2\sin\theta$, recuperando el factor clásico $dV = r^2\sin\theta\,dr\,d\theta\,d\phi$.

% =====================================================================
\manualsection[6. Diferenciales Exactos frente a Inexactos]{6. Física de las Formas Diferenciales: Diferenciales Exactos e Inexactos}

En termodinámica y mecánica elemental encontramos dos tipos de cantidades infinitesimales que suelen anotarse de forma similar pero obedecen leyes matemáticas radicalmente distintas:
\begin{itemize}
    \item \textbf{Diferenciales Exactos ($df$):} Corresponden a cambios infinitesimales de una \textbf{función de estado} o potencial $f(x, y, z)$. Por el Teorema Fundamental del Cálculo para integrales de línea, su valor depende exclusivamente de los puntos extremos:
    \[
    \int_{\mathcal{C}_A^B} df = f(B) - f(A), \qquad \oint_\mathcal{C} df = 0.
    \]
    Representan magnitudes conservativas, como la energía potencial electrostática o gravitatoria.
    \item \textbf{Diferenciales Inexactos ($\delta W, \delta Q$):} Corresponden a magnitudes que dependen de la \textbf{trayectoria específica} seguida en el espacio, como el trabajo mecánico realizado por una fuerza no conservativa o el calor en un ciclo termodinámico:
    \[
    W = \int_\gamma \delta W = \int_\gamma \vec{F}\cdot d\vec{r}.
    \]
    Para un diferencial inexacto, la circulación neta sobre un camino cerrado en general no se anula: $\oint_\mathcal{C} \delta W \neq 0$.
\end{itemize}

\manualsubsection{El Criterio de Integrabilidad de Schwarz}
¿Cómo determinamos si una forma diferencial dada $\omega = P(x,y)\,dx + Q(x,y)\,dy$ es exacta sin tener que adivinar previamente su función potencial?

\begin{theorembox}{Criterio de Integrabilidad de Schwarz}
Sea $\omega = P(x, y)\,dx + Q(x, y)\,dy$ una forma diferencial con coeficientes $\mathcal{C}^1$ en una región abierta y simplemente conexa. La forma $\omega$ es exacta ($\omega = df$) si y sólo si satisface la condición de simetría de derivadas cruzadas:
\[
\mathbf{\frac{\partial P}{\partial y} = \frac{\partial Q}{\partial x} \quad \Longleftrightarrow \quad \frac{\partial Q}{\partial x} - \frac{\partial P}{\partial y} = 0}.
\]
\end{theorembox}

\textbf{Demostración:}
\begin{enumerate}
    \item \textbf{Necesidad ($\implies$):} Si $\omega = df$, entonces por definición de diferencial total:
    \[
    df = \frac{\partial f}{\partial x} dx + \frac{\partial f}{\partial y} dy = P\,dx + Q\,dy \implies P = \frac{\partial f}{\partial x}, \quad Q = \frac{\partial f}{\partial y}.
    \]
    Derivando cruzadamente $P$ respecto a $y$ y $Q$ respecto a $x$:
    \[
    \frac{\partial P}{\partial y} = \frac{\partial^2 f}{\partial y \partial x}, \qquad \frac{\partial Q}{\partial x} = \frac{\partial^2 f}{\partial x \partial y}.
    \]
    Por el Teorema de Schwarz, las derivadas parciales de segundo orden de una función suave conmutan ($\frac{\partial^2 f}{\partial y \partial x} = \frac{\partial^2 f}{\partial x \partial y}$). Por lo tanto: $\frac{\partial P}{\partial y} = \frac{\partial Q}{\partial x}$.
    \item \textbf{Suficiencia ($\impliedby$):} En un dominio sin agujeros (estrellado o simplemente conexo), si $\partial_y P = \partial_x Q$, la función potencial se construye directamente mediante la integral de línea independiente de trayectoria:
    \[
    f(x, y) = \int_{(x_0, y_0)}^{(x, y)} \big(P(x', y')dx' + Q(x', y')dy'\big). \qquad \blacksquare
    \]
\end{enumerate}

\manualsubsection{La Topología del Dominio y el Campo de Vórtice}
¿Es suficiente que la condición local $\partial_y P = \partial_x Q$ se cumpla para asegurar que el trabajo en cualquier trayectoria cerrada sea cero? \textbf{La respuesta es no}, y la razón reside en la \textbf{topología del espacio}.

Consideremos el célebre campo vectorial en el plano perforado $\R^2 \setminus \{(0,0)\}$:
\[
\vec{F}(x, y) = \frac{-y\,\hat{i} + x\,\hat{j}}{x^2 + y^2} \iff \omega = \frac{-y\,dx + x\,dy}{x^2 + y^2}.
\]
Calculando las derivadas parciales cruzadas para cualquier punto fuera del origen:
\[
\frac{\partial P}{\partial y} = \frac{-(x^2 + y^2) - (-y)(2y)}{(x^2 + y^2)^2} = \frac{y^2 - x^2}{(x^2 + y^2)^2}, \qquad
\frac{\partial Q}{\partial x} = \frac{(x^2 + y^2) - x(2x)}{(x^2 + y^2)^2} = \frac{y^2 - x^2}{(x^2 + y^2)^2}.
\]
La condición de rotacional nulo $\partial_y P = \partial_x Q$ se cumple rigurosamente en todo el dominio. Sin embargo, si calculamos la circulación a lo largo de una circunferencia unitaria $\mathcal{C}$ parametrizada por $x(t) = \cos t$, $y(t) = \sin t$ con $t \in [0, 2\pi]$:
\[
\oint_\mathcal{C} \omega = \int_0^{2\pi} \frac{-\sin t(-\sin t\,dt) + \cos t(\cos t\,dt)}{\cos^2 t + \sin^2 t} = \int_0^{2\pi} 1\,dt = \mathbf{2\pi \neq 0}.
\]
¿Por qué falló el teorema? Porque localmente $\omega = d\theta$ (la diferencial del ángulo polar), pero el ángulo $\theta$ es una función multivaluada que da un salto de $2\pi$ al completar una vuelta. El agujero en el origen impide encoger la curva a un punto sin salirse del dominio.

\manualsubsection{El Salto a la Geometría Intrínseca}
Con este resultado culmina el **Módulo 00**:
\begin{enumerate}
    \item En la Semana 01 aprendimos a distinguir puntos de vectores.
    \item En la Semana 02 construimos el lenguaje indicial de Einstein y Levi-Civita.
    \item En la Semana 03 exploramos las transformaciones rígidas lineales $SO(3)$ y afines $SE(3)$.
    \item En esta Semana 04 descubrimos cómo la matriz Jacobiana $J$ y el tensor métrico $g_{ab} = J^T J$ permiten medir distancias reales en cualquier sistema de coordenadas no lineal.
\end{enumerate}
Este marco matemático nos deja en el umbral perfecto del **Módulo 01**, donde los vectores dejarán de verse como flechas euclidianas rígidas para revelarse como **derivaciones tangenciales**, y las formas diferenciales como **láminas geométricas** capaces de integrar la física sin depender de ningún sistema de referencia.

\end{document}
